\subsubsection{Edge-by-edge realization of the three nonadic extensions}
\label{subsubsec:tres-vueltas-ext-ch}

This subsection specifies the autonomous model generated by the TPK extension
relation entering the cardinal proof. The existence of an edge is not inferred
from a coincidence of types: its initial and final states, APP relation, fiber
transport, and unary update of the ledger are constructed. The construction
uses exclusively the two APP sheets, the TRIT selector, and the TPK
composition \(\mathsf{Sel}\)--\(\mathsf{Tra}\)--\(\operatorname{Upd}\).

\paragraph{Complete cells and three balanced reductions.}
In the chart \(D_9=\{1,\ldots,9\}\), we always retain the complete APP cell
\[
 a_{p,q}:=(p,q;R^\Sigma_{p,q},q^\Sigma_{p,q},
                 R^\Pi_{p,q},q^\Pi_{p,q}),
\]
with
\[
 p+q=R^\Sigma_{p,q}+9q^\Sigma_{p,q},\qquad
 pq=R^\Pi_{p,q}+9q^\Pi_{p,q}.
\]
The active sheet is recorded in a separate coordinate; \(a_{p,q}\) is not
split into a purported additive cell and a separate multiplicative cell. For
\(r\in D_9\), let
\begin{equation}
 a(r):=
 \begin{cases}
  a_{1,9},&r=1,\\
  a_{1,r-1},&2\le r\le9,
 \end{cases}
 \qquad
 \tau(r):=M_{\rm ph}(r),
 \qquad
 m(r):=\tau(r)\in\{-1,0,+1\}.
 \label{eq:celula-completa-fase-ch}
\end{equation}
The additive residue of \(a(r)\) supplies an APP representative of phase
here; this does not make \(\Sigma\) the universally active sheet. The dynamic
selector activates \(\Sigma\) for \(r\in\{1,4,7\}\), activates \(\Pi\) for
\(r\in\{2,5,8\}\), and preserves both sheets without moving the cursor for
\(r\in\{3,6,9\}\). In all three cases,
\(R^\Sigma,q^\Sigma,R^\Pi,q^\Pi\) are transported jointly. We have
\[
 (m(1),\ldots,m(9))=(1,-1,0,1,-1,0,1,-1,0).
\]
The three pairs ordered by the nonadic base are
\begin{equation}
 \mathfrak p_{+1}:=(1,4),\qquad
 \mathfrak p_{-1}:=(2,5),\qquad
 \mathfrak p_{0}:=(3,6).
 \label{eq:tres-pares-internos-ch}
\end{equation}
These pairs are the normalized section taking the first two occurrences of
each TRIT type under the shift \(r\mapsto r+3\); the third phase block
\(7,8,9\) retains the turn-closure control. Both members belong to the same
fiber of \(M_{\rm ph}\). Without introducing a target coefficient, the
balanced reduction from the TRIT chapter gives
\begin{equation}
 \begin{array}{c|c|c|c}
 \varepsilon&m(\mathfrak p_\varepsilon^{(1)})
                  +m(\mathfrak p_\varepsilon^{(2)})
   &r_3\bigl(2\varepsilon\bigr)&
   \chi(\varepsilon,\varepsilon)\\ \hline
 +1&1+1=(-1)+3\cdot1&-1&+1\\
 0&0+0=0+3\cdot0&0&0\\
 -1&(-1)+(-1)=1+3\cdot(-1)&+1&-1
 \end{array}
 \label{eq:tres-carry-derivados-ch}
\end{equation}
Therefore,
\begin{equation}
 \chi(\varepsilon,\varepsilon)
 =\frac{2\varepsilon-r_3(2\varepsilon)}3=\varepsilon.
 \label{eq:carry-no-parametrico-ch}
\end{equation}
The branch sign is thus the output of the balanced cocycle applied to two TRIT
emissions produced by APP. Each pair supplies one of the three carry outputs.
For clarity, we subsequently index the pair by
\(\varepsilon(\mathfrak p):=
\chi(m(\mathfrak p^{(1)}),m(\mathfrak p^{(2)}))\); the symbol
\(\varepsilon\) does not supply a free carry to \(\operatorname{Upd}^{\rm cal}\).

The phase architecture, balanced lift, and memory transport are realized here
through twenty-seven local incidences and their extension graphs. The
normalized section fixes representatives of the fibers of \(M_{\rm ph}\); the
covering to be proved uses the existence of these representatives and remains
invariant when they are replaced by another compatible section.

\paragraph{Typed data of the nine edges.}
Let \(K_{\rm Carr}:=\langle\kappa_{\rm Carr}\rangle_{\rm Ab}\) be the carry
module. The actions are typed by
\[
 H_{\rm Carr}:\mathscr P_{\rm TPK}^{\rm op}
   \longrightarrow\operatorname{End}(K_{\rm Carr}),\qquad
 Y_{\rm Carr}:\mathscr P_{\rm TPK}
   \longrightarrow\operatorname{End}(K_{\rm Carr}).
\]
In this generated model, take the trivial actions
\(H_{\rm Carr}(e)=Y_{\rm Carr}(e)=\operatorname{id}_{K_{\rm Carr}}\): the
active sheet and the route are retained as coordinates of the ledger itself
and do not act on the free carry summand. Also let
\((\mathsf M_d,\jmath_{d',d})\) be the graded direct system of memory. To
avoid confusing two representatives identified by the colimit, the state also
retains the degree as a coordinate, and the labeled carrier
\[
 \widetilde{\mathsf M}:=
 \bigsqcup_{d\ge1}\bigl(\{d\}\times\mathsf M_d\bigr),
 \qquad
 \operatorname{gr}_{\mathsf M}(d,\mu)=d.
\]
The representation in the limit memory is given by
\[
\begin{aligned}
 \mathsf M&:=\varinjlim_d\mathsf M_d,&
 \jmath_d&:\mathsf M_d\longrightarrow\mathsf M,\\
 u&:=\jmath_d(u_d),&
 \iota_{\mathsf M}:\mathbb Z&\longrightarrow\mathsf M,
 \qquad n\longmapsto nu .
\end{aligned}
\]
Here \(\jmath_d\) is the canonical colimit map, and the \(u_d\) form a
compatible collection. Its free component permits the reader \(\deg_u(nu)=n\).
The configuration, phase, memory, carry, route, and survival projections will
be typed on the calibrated carrier constructed recursively below; no ambient
space will be imported for them. For every constructed edge \(e\), the
calibrated component
\(\operatorname{Upd}^{\rm cal}_e:=\mathsf{Mem}^{\rm cal}_e\) will denote the
third map of the factorization; its law is fixed field by field below and is
not presupposed as the restriction of an ambient relation.

Fix the initial depth \(k_0=1\), a complete configuration
\(q_\star\in\mathcal Q_{\rm obs}\) of phase one, and
\[
 q_{n,r}:=F_{\rm obs}^{9n+r-1}(q_\star),\qquad
 q_{n,10}:=q_{n+1,1}.
\]
Because the observable calendar is periodic, its value \(q_{n,r}\) does not
by itself distinguish occurrences separated by twelve turns. The calibrated
model therefore uses the labeled lift
\[
 \widetilde q_{n,r}:=(q_{n,r},n,r),\qquad
 \widetilde q_{n,10}:=\widetilde q_{n+1,1},\qquad
 \pi_{\rm obs}(\widetilde q_{n,r})=q_{n,r}.
\]
The lifted dynamics
\(\widetilde F_{\rm obs}(\widetilde q_{n,r})
=\widetilde q_{n,r+1}\) projects to \(F_{\rm obs}\); the label does not alter the physical phase, but makes its
levels of occurrence distinct. The copy of the cell at each level is labeled by
\[
 a_{n,r}:=(n,r;a(r)),\qquad
 (n,r)^+:=
 \begin{cases}
  (n,r+1),&1\le r<9,\\
  (n+1,1),&r=9.
 \end{cases}
\]
The label only prevents two occurrences of the same cell from being
identified; the two Euclidean divisions and their four coordinates are those
of \(a(r)\). For turn \(n\ge0\) and phase \(r\in D_9\), the symbol
\[
 e_{\varepsilon,n,r}:\widetilde q_{n,r}
 \longrightarrow\widetilde q_{n,r+1},
 \qquad r^+:=1+(r\bmod9),
\]
occupies level \(d_{n,r}:=k_n+r-1\), where \(k_n:=k_0+9n\). Its local ledger
contains the complete cell \(a(r)\), the active sheet determined by
\(M_{\rm ph}(r)\), the type \(\tau(r)\), the membership mark in
\(\mathfrak p_\varepsilon\), and, upon reaching the second member of that
pair, the certified identity \eqref{eq:tres-carry-derivados-ch}. We define its
operators solely from that ledger:
\begin{align}
 \mathsf C_{\mathsf M}(e_{\varepsilon,n,r})
   &=\jmath_{d_{n,r}+1,d_{n,r}},&
 \kappa_{\mathsf M}(e_{\varepsilon,n,r})
   &=\mathbf1_{\{r=9\}}u_{d_{n,r}+1},
 \label{eq:memoria-arista-ch}\\
       \operatorname{Carr}(e_{\varepsilon,n,r})
   &=d_{\varepsilon,r}\kappa_{\rm Carr},&
 d_{\varepsilon,r}
   &:=\mathbf1_{\{r=\mathfrak p_\varepsilon^{(2)}\}}
       \chi\!\left(m(\mathfrak p_\varepsilon^{(1)}),
                    m(\mathfrak p_\varepsilon^{(2)})\right).
 \label{eq:carry-arista-derivado-ch}
\end{align}
For identities and every admissible composition, also define
\begin{align}
 \mathsf C_{\mathsf M}(\operatorname{id}_{\widetilde q};d)
   &:=\operatorname{id}_{\mathsf M_d},&
 \kappa_{\mathsf M}(\operatorname{id}_{\widetilde q};d)&:=0,
 \label{eq:memoria-identidad-ch}\\
 \mathsf C_{\mathsf M}(e_t\cdots e_1)
   &:=\mathsf C_{\mathsf M}(e_t)\circ\cdots\circ
      \mathsf C_{\mathsf M}(e_1),&
 \kappa_{\mathsf M}(e_2e_1)
   &:=\mathsf C_{\mathsf M}(e_2)
        \bigl(\kappa_{\mathsf M}(e_1)\bigr)
      +\kappa_{\mathsf M}(e_2).
 \label{eq:cociclo-memoria-calibrado-ch}
\end{align}
The identity carries degree \(d\) because the same observable configuration
may recur at different depths. Iterated, the final equality defines
\(\kappa_{\mathsf M}(e_t\cdots e_1)\) for every length and is the cocycle law
of the direct system; the inclusions satisfy
\(\jmath_{d+2,d+1}\jmath_{d+1,d}=\jmath_{d+2,d}\).
In addition, set \(s(e_{\varepsilon,n,r})=\top\). The following twenty-seven
entries certify the coordinates distinguishing the three extensions; the
universal transport fields are specified immediately afterward. The symbols
\(I\) and \(II\) designate, respectively, the first retained emission and the
second emission performing the reduction; a dash denotes a zero carry
increment, not the absence of an edge.
\begin{center}
\small
\setlength{\tabcolsep}{3.2pt}
\begin{tabular}{c c cc cc cc c}
\toprule
\(r\)&\(M_{\rm ph}(r)\)&\multicolumn{2}{c}{\(\mathfrak p_{+1}\)}&
\multicolumn{2}{c}{\(\mathfrak p_0\)}&
\multicolumn{2}{c}{\(\mathfrak p_{-1}\)}&
\(\Delta\operatorname{mem}\)\\
& &mark&\(d_{+1,r}\)&mark&\(d_{0,r}\)&mark&\(d_{-1,r}\)&\\
\midrule
1&\(+1\)&I&0&--&0&--&0&0\\
2&\(-1\)&--&0&--&0&I&0&0\\
3&0&--&0&I&0&--&0&0\\
4&\(+1\)&II&\(+1\)&--&0&--&0&0\\
5&\(-1\)&--&0&--&0&II&\(-1\)&0\\
6&0&--&0&II&0&--&0&0\\
7&\(+1\)&--&0&--&0&--&0&0\\
8&\(-1\)&--&0&--&0&--&0&0\\
9&0&--&0&--&0&--&0&\(u\)\\
\bottomrule
\end{tabular}
\end{center}
Thus, \(d_{\varepsilon,r}\) is computed from two outputs of the selector and
is distinct from the cell quotients \(q^\Sigma,q^\Pi\). The ledger also
preserves the marked pair and the balanced residue; in particular, the zero
branch is not confused with an absence of event.

The carry composition law used by the TPK category is
\begin{equation}
 \operatorname{Carr}(\operatorname{id}_{\widetilde q})=0,\qquad
 \operatorname{Carr}(e_2e_1)
 =H_{\rm Carr}(e_1)\bigl(\operatorname{Carr}(e_2)\bigr)
  +Y_{\rm Carr}(e_2)\bigl(\operatorname{Carr}(e_1)\bigr).
 \label{eq:ley-carry-correcta-ch}
\end{equation}
On identities and paths, they are extended functorially:
\[
\begin{aligned}
 H_{\rm Carr}(\operatorname{id})
   &=Y_{\rm Carr}(\operatorname{id})=\operatorname{id},\\
 H_{\rm Carr}(e_t\cdots e_1)
   &=H_{\rm Carr}(e_1)\circ\cdots\circ H_{\rm Carr}(e_t),\\
 Y_{\rm Carr}(e_t\cdots e_1)
   &=Y_{\rm Carr}(e_t)\circ\cdots\circ Y_{\rm Carr}(e_1).
\end{aligned}
\]
In this model all these actions remain the identity, but their typing makes
explicit which endomorphism acts on each summand of
\eqref{eq:ley-carry-correcta-ch}. On the calibrated edges,
\(H_{\rm Carr}=Y_{\rm Carr}=\operatorname{id}\), so composition adds the
increments already produced. Equation~\eqref{eq:ley-carry-correcta-ch}, not an
identification between orientation and quotient, governs the accumulated
ledger.

\paragraph{Independent construction of the transition tuples.}
A fiber state is written, in the order of its relevant fields, as
\[
 x=(a;q;\tau,o,h;b_{\rm ref};
       R^\Sigma,q^\Sigma,R^\Pi,q^\Pi;c;
       \operatorname{Sig};C,\partial C;p_{\APP},p_{\TRIT};
       \lambda;\operatorname{gr}_{\mathsf M};\operatorname{mem};
       s;\operatorname{Led}).
\]
The update notation \(x[f\leftarrow v]\) changes only the indicated field.
For \(h\in\{(\Sigma,\Pi),(\Pi,\Sigma)\}\), define
\[
 h_r(h)=
 \begin{cases}
  (\Sigma,\Pi),&r\in\{1,4,7\},\\
  (\Pi,\Sigma),&r\in\{2,5,8\},\\
 h,&r\in\{3,6,9\}.
 \end{cases}
\]
The complete sheet-and-arithmetic coordinate is abbreviated as
\[
 \widehat h(x):=
 \bigl(h(x),R^\Sigma(x),q^\Sigma(x),R^\Pi(x),q^\Pi(x)\bigr).
\]
Let \(\widehat{\mathcal H}^{\rm cal}_{\widetilde q_{n,r}}\) be the set of
quintuples
\[
 \bigl(h,R^\Sigma_{a_{n,r}},q^\Sigma_{a_{n,r}},
          R^\Pi_{a_{n,r}},q^\Pi_{a_{n,r}}\bigr),
 \qquad
 h\in\{(\Sigma,\Pi),(\Pi,\Sigma)\},
\]
whose four arithmetic data are the two Euclidean divisions of the labeled APP
cell \(a_{n,r}\). This direct definition does not presuppose a state in a
carrier not yet constructed. Let
\(\mathsf{Hist}^{\rm cal}_{\widetilde q}\) be the set of admissible finite
paths in the graph of labeled edges ending at \(\widetilde q\), including the
identity, and let \(\mathsf B^{\rm cal}_{\widetilde q}\) be the set of their
ordered endpoint sequences. Define
\[
 \mathsf S^{\rm cal}_{\widetilde q}:=
 \widehat{\mathcal H}^{\rm cal}_{\widetilde q}
 \times\{o_{\rightarrow},o_{\circ},o_{\leftarrow}\}
 \times K_{\rm Carr}\times
 \mathsf{Hist}^{\rm cal}_{\widetilde q}\times
 \mathsf B^{\rm cal}_{\widetilde q}\times\mathsf M .
\]
The calibrated signature is not postulated: it is the typed assembly map
\[
 \operatorname{Sig}^{\rm cal}_{\widetilde q}:
 \widehat{\mathcal H}^{\rm cal}_{\widetilde q}
 \times\{o_{\rightarrow},o_{\circ},o_{\leftarrow}\}
 \times K_{\rm Carr}\times
 \mathsf{Hist}^{\rm cal}_{\widetilde q}\times
 \mathsf B^{\rm cal}_{\widetilde q}\times\mathsf M
 \longrightarrow\mathsf S^{\rm cal}_{\widetilde q},
\quad
 (\widehat h,o,c,\lambda,\partial C,\mu)
 \longmapsto(\widehat h,o,c,\lambda,\partial C,\mu).
\]
Likewise, \(\pi_{\mathcal F}\) denotes the projection that removes only the
first two coordinates \((a;q)\) of the state tuple; it preserves
\(\tau,o,h,b_{\rm ref}\), the four Euclidean data, carry, signature, cylinder,
boundary, prefixes, route, memory degree, memory, survival, and ledger. The
corresponding orientation is fixed by
\(o(+1)=o_{\rightarrow}\), \(o(0)=o_{\circ}\), and
\(o(-1)=o_{\leftarrow}\).
The calibrated APP fiber over \(\widetilde q_{n,r}\) contains the already
defined labeled copy \(a_{n,r}\). The initial state is fixed without free
coordinates:
\[
 \widehat h_\star:=
 \bigl((\Sigma,\Pi),R^\Sigma_{1,9},q^\Sigma_{1,9},
       R^\Pi_{1,9},q^\Pi_{1,9}\bigr).
\]
\begin{equation}
\begin{aligned}
 x_\star:=\bigl(&a_{0,1};\widetilde q_{0,1};
 +1,o_{\rightarrow},(\Sigma,\Pi);\bot;
 R^\Sigma_{1,9},q^\Sigma_{1,9},R^\Pi_{1,9},q^\Pi_{1,9};0;\\
 &\operatorname{Sig}^{\rm cal}_{\widetilde q_{0,1}}
       (\widehat h_\star,o_{\rightarrow},0,
        \operatorname{id}_{\widetilde q_{0,1}},\varnothing,
        \jmath_{k_0}(0_{\mathsf M_{k_0}}));
 \varnothing,\varnothing;
 (a_{0,1}),(+1);\operatorname{id}_{\widetilde q_{0,1}};
 k_0;0_{\mathsf M_{k_0}};\top;\varnothing\bigr).
\end{aligned}
\label{eq:estado-inicial-calibrado-ch}
\end{equation}
The tuple contains no indeterminate coordinates and satisfies
\(s(x_\star)=\top\). Let \(G_0:=\{x_\star\}\). Assuming \(G_n\) has been
constructed, fix \(x\in G_n\), \(\varepsilon\in\{-1,0,+1\}\), and
\(x_{\varepsilon,0}(x):=x\). For \(r=1,\ldots,9\), first construct the
elementary ledger
\begin{equation}
\begin{aligned}
 \ell_{\varepsilon,n,r}:=\bigl(&e_{\varepsilon,n,r},a_{n,r},a_{(n,r)^+},
 \widetilde q_{n,r},\widetilde q_{n,r+1},M_{\rm ph}(r),\\
 &h_r(h(x_{\varepsilon,r-1}(x))),o(M_{\rm ph}(r)),\mathfrak p_\varepsilon,
 d_{\varepsilon,r},\\
 &\mathsf C_{\mathsf M}(e_{\varepsilon,n,r}),
 \kappa_{\mathsf M}(e_{\varepsilon,n,r}),
 \operatorname{Carr}(e_{\varepsilon,n,r})\bigr),
\end{aligned}
 \label{eq:registro-elemental-ext-ch}
\end{equation}
and set
\(\operatorname{Led}^{\rm new}_{\varepsilon,n,r}:=
\operatorname{Led}(x_{\varepsilon,r-1}(x))^\frown
\ell_{\varepsilon,n,r}\). Then construct the tuple
\begin{equation}
 y_{\varepsilon,r}(x):=
 x_{\varepsilon,r-1}(x)
 [\tau\leftarrow M_{\rm ph}(r),\
  h\leftarrow h_r(h(x_{\varepsilon,r-1}(x))),\
  o\leftarrow o(M_{\rm ph}(r)),\
  b_{\rm ref}\leftarrow\mathfrak p_\varepsilon],
 \label{eq:tupla-sel-cal-ch}
\end{equation}
where selection leaves the cell, quotients, carry, route, memory, cylinder,
boundary, survival, and ledger unchanged; the coordinate \(b_{\rm ref}\)
makes the chosen balanced reduction explicit. Membership in an ambient
relation is not presupposed here.

Next, define \(z_{\varepsilon,r}(x)\) field by field: its configuration is
\(\widetilde q_{n,r+1}\); its complete cell is \(a_{(n,r)^+}\), with the four
residues and quotients recomputed by the two Euclidean divisions; it preserves
the selected sheet and orientation; and it performs the following six updates,
in which only the common argument \(x\) is omitted:
\begin{equation}
\begin{aligned}
 C(z_{\varepsilon,r})&:=C(y_{\varepsilon,r})^\frown
       e_{\varepsilon,n,r},\\
 \partial C(z_{\varepsilon,r})&:=\partial C(y_{\varepsilon,r})^\frown
       (\widetilde q_{n,r},\widetilde q_{n,r+1}),\\
 p_{\APP}(z_{\varepsilon,r})&:=p_{\APP}(y_{\varepsilon,r})^\frown
       a_{(n,r)^+},\\
 p_{\TRIT}(z_{\varepsilon,r})&:=p_{\TRIT}(y_{\varepsilon,r})^\frown
       \tau(r),\\
 \operatorname{Led}(z_{\varepsilon,r})&:=
       \operatorname{Led}^{\rm new}_{\varepsilon,n,r},\\
 \lambda(z_{\varepsilon,r})&:=e_{\varepsilon,n,r}
       \lambda(y_{\varepsilon,r}),
\end{aligned}
 \label{eq:tupla-tra-cal-ch}
\end{equation}
and still retains the fields
\(c,\operatorname{gr}_{\mathsf M},\operatorname{mem},s\) of
\(y_{\varepsilon,r}(x)\). Its signature, by contrast, is retyped at the target
configuration through the transported presignature
\begin{equation}
\begin{aligned}
 \operatorname{Sig}(z_{\varepsilon,r}(x))
  :=\operatorname{Sig}^{\rm cal}_{\widetilde q_{n,r+1}}
       \Bigl(&\widehat h(z_{\varepsilon,r}(x)),
             o(z_{\varepsilon,r}(x)),c(y_{\varepsilon,r}(x)),
             \lambda(z_{\varepsilon,r}(x)),\\[-2pt]
             &\partial C(z_{\varepsilon,r}(x)),
             \jmath_{d_{n,r}}
                (\operatorname{mem}_{d_{n,r}}
                   (y_{\varepsilon,r}(x)))\Bigr).
\end{aligned}
\label{eq:prefirma-transporte-cal-ch}
\end{equation}
Thus, and only in \(\mathsf{Tra}^{\rm cal}\), the edge, its endpoints, the
mark \((\mathfrak p_\varepsilon,r,d_{\varepsilon,r})\), and its cocycles are
added to the incidence ledger.

The initial state satisfies
\[
 c(x_\star)=0=\operatorname{Carr}
   (\operatorname{id}_{\widetilde q_{0,1}})
   =\operatorname{Carr}(\lambda(x_\star)).
\]
In the complete states \(x_{\varepsilon,r-1}(x)\), and also after the
selection \(y_{\varepsilon,r}(x)\), the typed invariant
\(c(v)=\operatorname{Carr}(\lambda(v))\) is preserved, since
\(\mathsf{Sel}^{\rm cal}\) modifies neither field. Transport advances the
route but leaves the carry pending; hence the pre-update tuple satisfies
exactly
\[
 c(z_{\varepsilon,r}(x))=c(y_{\varepsilon,r}(x))
 =\operatorname{Carr}(\lambda(y_{\varepsilon,r}(x))).
\]
The map \(\operatorname{Upd}^{\rm cal}\) then restores
\(c(x_{\varepsilon,r}(x))=
\operatorname{Carr}(\lambda(x_{\varepsilon,r}(x)))\) by means of the
composition law, without attributing that invariant to the intermediate state \(z\).

Finally, apply the unary function
\(\operatorname{Upd}^{\rm cal}_{e_{\varepsilon,n,r}}
=\mathsf{Mem}^{\rm cal}_{e_{\varepsilon,n,r}}\) to that tuple. Its result
\(x_{\varepsilon,r}(x)\) is given by
\begin{equation}
\begin{aligned}
 c\bigl(x_{\varepsilon,r}(x)\bigr)
   &=\operatorname{Carr}\bigl(\lambda(z_{\varepsilon,r}(x))\bigr)\\
   &=H_{\rm Carr}(\lambda(y_{\varepsilon,r}(x)))
       \bigl(\operatorname{Carr}(e_{\varepsilon,n,r})\bigr)
     +Y_{\rm Carr}(e_{\varepsilon,n,r})
       \bigl(c(y_{\varepsilon,r}(x))\bigr),\\
 \lambda\bigl(x_{\varepsilon,r}(x)\bigr)
   &=\lambda(z_{\varepsilon,r}(x)),\\
 \operatorname{gr}_{\mathsf M}
      \bigl(x_{\varepsilon,r}(x)\bigr)
   &=d_{n,r}+1,\\
 \operatorname{mem}_{d_{n,r}+1}
       \bigl(x_{\varepsilon,r}(x)\bigr)
   &=\mathsf C_{\mathsf M}(e_{\varepsilon,n,r})
       \operatorname{mem}_{d_{n,r}}(z_{\varepsilon,r}(x))
     +\kappa_{\mathsf M}(e_{\varepsilon,n,r}),\\
 s\bigl(x_{\varepsilon,r}(x)\bigr)
   &=s(e_{\varepsilon,n,r})\wedge s(z_{\varepsilon,r}(x)),\\
 \operatorname{Sig}\bigl(x_{\varepsilon,r}(x)\bigr)
   &=\operatorname{Sig}^{\rm cal}_{\widetilde q_{n,r+1}}
       \Bigl(\widehat h(z_{\varepsilon,r}(x)),
             o(z_{\varepsilon,r}(x)),
             c(x_{\varepsilon,r}(x)),
             \lambda(z_{\varepsilon,r}(x)),
             \partial C(z_{\varepsilon,r}(x)),\\[-2pt]
   &\hspace{112pt}
             \jmath_{d_{n,r}+1}
              \bigl(\operatorname{mem}_{d_{n,r}+1}
                 (x_{\varepsilon,r}(x))\bigr)\Bigr).
\end{aligned}
\label{eq:upd-unaria-tres-vueltas-ch}
\end{equation}
All fields not enumerated in \eqref{eq:upd-unaria-tres-vueltas-ch} coincide
with those of \(z_{\varepsilon,r}(x)\). In particular, the signature is
computed from the sheet, orientation, boundary, and route already fixed by
\(\mathsf{Sel}^{\rm cal}\) and \(\mathsf{Tra}^{\rm cal}\), and from the carry
and memory computed in the preceding lines; it is not circularly defined from
the final state. Neither \(\varepsilon\), nor \(d_{\varepsilon,r}\), nor
\(u\) is an additional argument of \(\operatorname{Upd}^{\rm cal}\): the
function reads all three data from the already constructed edge.

\paragraph{Simultaneous induction of the levels.}
Once the preceding nine stages have been completed for an already constructed
\(G_n\), define immediately
\begin{equation}
 \gamma_{\varepsilon,k_n}(x):=x_{\varepsilon,9}(x),\qquad
 G_{n+1}:=\{\gamma_{\varepsilon,k_n}(x):
       x\in G_n,\ \varepsilon\in\TRIT\}.
 \label{eq:injerto-total-tres-vueltas-ch}
\end{equation}
\label{eq:tres-vueltas-tpk-ch}
The base \(G_0=\{x_\star\}\), the intermediate tuples, and
\eqref{eq:injerto-total-tres-vueltas-ch} constitute a single definition by
recursion on \(n\). Therefore, no subsequent carrier is used to assume the
existence of \(G_n\): the next level is introduced only after constructing all
its edges from the preceding level.

\paragraph{TPK extension relation and effective membership.}
Let \(\mathcal E^{\rm cal}_{\widetilde q}\) be the least labeled collection
containing \(x_\star\) and, at each step of the preceding recursion, the four
tuples
\(x_{\varepsilon,r-1}(x),y_{\varepsilon,r}(x),
z_{\varepsilon,r}(x),x_{\varepsilon,r}(x)\) in their source and target
configurations. Also let
\(\mathcal A^{\rm cal}_{\widetilde q_{n,r}}:=\{a_{n,r}\}\), and set
\[
 \widetilde{\mathcal Q}_{\rm obs}:=
   \{\widetilde q_{n,r}:n\ge0,\ 1\le r\le9\},
 \qquad
 \mathcal E^{\rm cal}:=\bigsqcup_{\widetilde q}
     \mathcal E^{\rm cal}_{\widetilde q},
 \qquad
 \mathcal F^{\rm cal}_{\widetilde q}:=
     \pi_{\mathcal F}(\mathcal E^{\rm cal}_{\widetilde q}).
\]
For each depth \(d\), let
\(\mathcal E^{\rm cal}_d:=
\operatorname{gr}_{\mathsf M}^{-1}(d)\subset\mathcal E^{\rm cal}\). The
explicit degree coordinate, rather than the mere membership of a
representative in some image of the direct system, makes these carriers
disjoint. On them, the projections are typed as
\[
\begin{aligned}
 \operatorname{cfg}&:\mathcal E^{\rm cal}\longrightarrow
     \widetilde{\mathcal Q}_{\rm obs},\\
 g(x)&:=\operatorname{phase}\!
       \bigl(\pi_{\rm obs}(\operatorname{cfg}(x))\bigr)\in C_9,\\
 c&:\mathcal E^{\rm cal}\longrightarrow K_{\rm Carr},\\
 s&:\mathcal E^{\rm cal}\longrightarrow\{\bot,\top\},\\
 \lambda&:\mathcal E^{\rm cal}_{\widetilde q}
       \longrightarrow\mathsf{Hist}^{\rm cal}_{\widetilde q},\\
 \operatorname{gr}_{\mathsf M}&:\mathcal E^{\rm cal}
       \longrightarrow\mathbb N_{\ge1},\\
 \operatorname{mem}_d&:\mathcal E^{\rm cal}_d
       \longrightarrow\mathsf M_d,\\
 \operatorname{mem}(x)&:=
       \jmath_d\bigl(\operatorname{mem}_d(x)\bigr)\in\mathsf M
       \quad(x\text{ of depth }d).
\end{aligned}
\]
Here \(\operatorname{cfg}(x)\) is the second coordinate of the tuple, and
\(\operatorname{phase}\) reads the phase of the observable calendar. For the
three carrier classes, fix the diagonals
\[
\begin{aligned}
 \Delta_{\APP,\widetilde q}
   &:=\{(a,a):a\in\mathcal A^{\rm cal}_{\widetilde q}\},\\
 \Delta_{{\rm fib},\widetilde q}
   &:=\{(f,f):f\in\mathcal F^{\rm cal}_{\widetilde q}\},\\
 \Delta_{{\rm enr},\widetilde q}
   &:=\{(x,x):x\in\mathcal E^{\rm cal}_{\widetilde q}\}.
\end{aligned}
\]
In particular,
\[
 a_{0,1}\in\mathcal A^{\rm cal}_{\widetilde q_{0,1}},
 \qquad
 x_\star\in\mathcal E^{\rm cal}_{\widetilde q_{0,1}}.
\]
On these carriers, the following relations are defined by their complete
graphs, and not through necessary implications:
\begin{align}
 \mathsf{Sel}^{\rm cal}_{e_{\varepsilon,n,r}}
   &:=\{(x_{\varepsilon,r-1}(x),y_{\varepsilon,r}(x)):
          x\in G_n\},\\
 \mathsf{Tra}^{\rm cal}_{e_{\varepsilon,n,r}}
   &:=\{(y_{\varepsilon,r}(x),z_{\varepsilon,r}(x)):
          x\in G_n\},\\
 \operatorname{Ext}^{\APP,\rm cal}_{e_{\varepsilon,n,r}}
   &:=\{(a_{n,r},a_{(n,r)^+})\},
 \label{eq:ext-app-calibrada-ch}\\
 \operatorname{Ext}^{\rm fib,cal}_{e_{\varepsilon,n,r}}
   &:=\{(\pi_{\mathcal F}x_{\varepsilon,r-1}(x),
          \pi_{\mathcal F}x_{\varepsilon,r}(x)):x\in G_n\}.
 \label{eq:ext-fibra-calibrada-ch}
\end{align}
The synchronized enriched extension is
\begin{equation}
 \operatorname{Ext}^{\rm enr,cal}_{e_{\varepsilon,n,r}}
 :=\{(x_{\varepsilon,r-1}(x),x_{\varepsilon,r}(x)):
       x\in G_n\}.
 \label{eq:ext-enriquecida-calibrada-ch}
\end{equation}
The third map is defined by
\begin{equation}
 \operatorname{graph}
 \bigl(\operatorname{Upd}^{\rm cal}_{e_{\varepsilon,n,r}}\bigr)
 :=\{(z_{\varepsilon,r}(x),x_{\varepsilon,r}(x)):
       x\in G_n\},
 \label{eq:grafo-upd-calibrado-ch}
\end{equation}
where the second component is exactly the one in
\eqref{eq:upd-unaria-tres-vueltas-ch}. Consequently,
\begin{equation}
 \mathcal U^{(1),\rm cal}_{e_{\varepsilon,n,r}}
 :=\operatorname{Upd}^{\rm cal}_{e_{\varepsilon,n,r}}\circ
   \mathsf{Tra}^{\rm cal}_{e_{\varepsilon,n,r}}\circ
   \mathsf{Sel}^{\rm cal}_{e_{\varepsilon,n,r}}
 =\{(x_{\varepsilon,r-1}(x),x_{\varepsilon,r}(x)):
      x\in G_n\}.
 \label{eq:transicion-cal-en-tpk-ch}
\end{equation}
Each operator is indexed by the edge \(e_{\varepsilon,n,r}\) produced by the
reduction \(\mathfrak p_\varepsilon\), rather than by a trit supplied from
outside. The coordinate \(b_{\rm ref}\) also distinguishes their codomains.
Therefore,
\(\mathsf{Sel}^{\rm cal}_{e}\),
\(\mathsf{Tra}^{\rm cal}\) and \(\operatorname{Upd}^{\rm cal}\) are functions
on their respective domains, whereas the union
\(\bigcup_{\varepsilon\in\TRIT}
\mathcal U^{(1),\rm cal}_{e_{\varepsilon,n,r}}\) deliberately preserves the
three outputs of the trisection.

For \(\bullet\in\{\APP,{\rm fib},{\rm enr}\}\), identities and compositions
are not left implicit. Define
\begin{align}
 \operatorname{Ext}^{\bullet,\rm cal}_{\operatorname{id}_{\widetilde q}}
   &:=\Delta_{\bullet,\widetilde q},\\
 \operatorname{Ext}^{\bullet,\rm cal}_{e_t\cdots e_1}
   &:=\operatorname{Ext}^{\bullet,\rm cal}_{e_t}\circ\cdots\circ
      \operatorname{Ext}^{\bullet,\rm cal}_{e_1}
 \label{eq:ext-palabra-calibrada-ch}
\end{align}
for every admissible edge word. The two APP and fiber path closures are formed
separately,
\begin{equation}
 \operatorname{Path}(\operatorname{Ext}^{\APP,\rm cal}),\qquad
 \operatorname{Path}(\operatorname{Ext}^{\rm fib,cal}),
 \label{eq:dos-clausuras-ext-ch}
\end{equation}
and \(\mathsf{Tra}^{\rm cal}\) synchronizes them edge by edge.

Finally, the levels of complete states and their truncations are fixed by the
recursion itself:
\begin{align}
 H^{\rm cal}_{d_{n,r}}
   &:=\{x_{\varepsilon,r-1}(x):
          x\in G_n,\ \varepsilon\in\TRIT\},\\
 H^{\rm cal}_{d_{n,r}+1}
   &:=\{x_{\varepsilon,r}(x):
          x\in G_n,\ \varepsilon\in\TRIT\},\\
 \rho^{\rm cal}_{d_{n,r}+1,d_{n,r}}
   \bigl(x_{\varepsilon,r}(x)\bigr)
   &:=x_{\varepsilon,r-1}(x).
 \label{eq:truncamiento-calibrado-definido-ch}
\end{align}
The final ledger preserves \((n,r,\mathfrak p_\varepsilon)\), so the
predecessor of the left-hand side is unique and \(\rho^{\rm cal}\) is
well-defined. For \(b>a\), set
\(\rho^{\rm cal}_{b,a}:=\rho^{\rm cal}_{a+1,a}\circ\cdots\circ
\rho^{\rm cal}_{b,b-1}\). This map recovers the complete preceding state,
including its signature; it does not purport to invert an overwritten
coordinate without consulting the ledger.

\begin{proposition}[Autonomous realization of the TPK through explicit transitions]
\label{prop:modelo-tpk-calibrado-ch}
The collection
\[
\begin{aligned}
 \mathfrak T_{\rm TPK}^{\rm cal}:=\bigl(
 &(\widetilde q_{n,r})_{n,r},\pi_{\rm obs},
   \widetilde F_{\rm obs},M_{\rm ph},\\
 &(\mathcal A_{\widetilde q}^{\rm cal})_{\widetilde q},
   (\mathcal E_{\widetilde q}^{\rm cal})_{\widetilde q},
   (\mathcal F_{\widetilde q}^{\rm cal})_{\widetilde q},
   \operatorname{Ext}^{\APP,\rm cal},
   \operatorname{Ext}^{\rm fib,cal},
   \operatorname{Ext}^{\rm enr,cal},\rho^{\rm cal},\\
 &\mathsf{Sel}^{\rm cal},\mathsf{Tra}^{\rm cal},
   \operatorname{Upd}^{\rm cal},\mathsf C_{\mathsf M},
   \kappa_{\mathsf M},\operatorname{gr}_{\mathsf M},
   H_{\rm Carr},Y_{\rm Carr},\\
 &\operatorname{Sig}^{\rm cal},\pi_{\mathcal F}\bigr).
\end{aligned}
\]
is an enriched model of the TPK. In particular, every pair on the right-hand
side of \eqref{eq:transicion-cal-en-tpk-ch} is an effective TPK edge, not a
consequence inferred from a merely necessary condition.
\end{proposition}

\begin{proof}
The labels \((n,r)\) make the carrier copies disjoint and unambiguously fix
each source and target. The relation \(\operatorname{Ext}^{\APP,\rm cal}\)
transports the complete cell and preserves the two Euclidean identities; the
relation \(\operatorname{Ext}^{\rm fib,cal}\) relates the fibers of the
preceding and subsequent complete states: it transports sheet, orientation,
prefix, cylinder, boundary, and ledger, and incorporates the updated carry,
memory, and signature. Selection does not modify the ledger, and transport
concatenates exactly one entry and produces the target presignature
\eqref{eq:prefirma-transporte-cal-ch}. The update recomputes the signature by
means of the typed map \(\operatorname{Sig}^{\rm cal}_{\widetilde q}\) from
the already determined sheet, orientation, carry, route, boundary, and memory.
The map \(\operatorname{Upd}^{\rm cal}\) is unary because the edge contained
in \(z_{\varepsilon,r}(x)\) determines its cocycles.

The identities are the diagonals, and the word extensions are the compositions
in \eqref{eq:ext-palabra-calibrada-ch}. Concatenation of prefixes and ledgers
is associative; memory satisfies the cocycle law, and carry satisfies
\eqref{eq:ley-carry-correcta-ch}. Therefore, the two parenthesizations of three
edges produce the same state. Formula
\eqref{eq:truncamiento-calibrado-definido-ch} defines truncation on complete
states, and uniqueness of the final ledger proves that it recovers the source
in every field. These are the identity, source, target, composition, memory,
carry, and truncation laws of the enriched TPK. In particular, the extension
relations are determined as graphs of effective maps rather than as
consequences of merely necessary conditions.
\end{proof}

Each \(\gamma_{\varepsilon,k_n}\) comprises exactly nine applications
of \(\operatorname{Upd}^{\rm cal}_{e_{\varepsilon,n,r}}\circ
\mathsf{Tra}^{\rm cal}_{e_{\varepsilon,n,r}}\circ
\mathsf{Sel}^{\rm cal}_{e_{\varepsilon,n,r}}\). Survival will be proved by
means of a constructed tail and will not be used to define \(G_n\).

\begin{lemma}[Effective internal trisection of one turn]
\label{lem:elevacion-catalogal-nonadica}
For every \(n\ge0\), every \(x\in G_n\), and every
\(\varepsilon\in\{-1,0,+1\}\), the expression
\(\gamma_{\varepsilon,k_n}(x)\) is defined by nine TPK transitions, with its
synchronized relations
\(\operatorname{Ext}^{\APP,\rm cal}\) and
\(\operatorname{Ext}^{\rm fib,cal}\), and satisfies
\begin{align}
 \rho^{\rm cal}_{k_n+9,k_n}\gamma_{\varepsilon,k_n}(x)&=x,
 \label{eq:truncamiento-injerto-ch}\\
 g\bigl(\gamma_{\varepsilon,k_n}(x)\bigr)
   &=g(x),&
 \jmath_{k_n+9}\!\left(\operatorname{mem}_{k_n+9}
        (\gamma_{\varepsilon,k_n}(x))\right)
   &=\jmath_{k_n}\!\left(\operatorname{mem}_{k_n}(x)\right)+u,
 \label{eq:fase-memoria-injerto-ch}\\
 c\bigl(\gamma_{\varepsilon,k_n}(x)\bigr)
   &=c(x)+\varepsilon\kappa_{\rm Carr}.&&
 \label{eq:carry-injerto-ch}
\end{align}
\label{eq:tres-vueltas-propiedades-ch}
\label{eq:tres-vueltas-acarreo-ch}
The ordered ledger of the nine edges determines \(\varepsilon\) uniquely.
The three images are therefore disjoint. Each endpoint belongs to
\(G_{n+1}\) and again admits the three maps
\(\gamma_{-1,k_{n+1}},\gamma_{0,k_{n+1}},
\gamma_{+1,k_{n+1}}\).
\end{lemma}

\begin{proof}
Formulas \eqref{eq:celula-completa-fase-ch} traverse the nine phases once and
return the copy \(a_{n+1,1}\) of \(a(1)\). At each step, the APP relation
transports the complete sextuple; \(\mathsf{Sel}^{\rm cal}\) fixes the TRIT
type from the phase mark \(r\) preserved in the state; and
\(\mathsf{Tra}^{\rm cal}\) adds exactly one symbol to the prefix, one segment
to the cylinder, and its ordered pair of endpoints to the boundary. By
induction on \(r\), the explicit tuples \(y_{\varepsilon,r}(x)\) and
\(z_{\varepsilon,r}(x)\) satisfy the predicates of the indicated relations,
and the unary update produces the state at the next level. The identity
\(\rho^{\rm cal}_{d_{n,r}+1,d_{n,r}}
(x_{\varepsilon,r}(x))=x_{\varepsilon,r-1}(x)\) is the definition
\eqref{eq:truncamiento-calibrado-definido-ch}; its well-definedness follows
from uniqueness of the final ledger. The composition of the nine truncation
maps recovers \(x\), proving \eqref{eq:truncamiento-injerto-ch} without
separately inverting the overwritten fields.

The phase \(9\to1\) occurs exactly once. By
\eqref{eq:memoria-arista-ch} and the cocycle law, the linear memory component
transports the preceding memory, and its translational component adds exactly
\(u\) in the direct limit. The compatibility
\(\jmath_{d+1}\jmath_{d+1,d}=\jmath_d\) proves
\eqref{eq:fase-memoria-injerto-ch}; the return concerns the phase, not the
enriched state.

By \eqref{eq:carry-arista-derivado-ch}, eight edges have zero increment, and
the edge completing \(\mathfrak p_\varepsilon\) has the computed increment
\(\chi(\varepsilon,\varepsilon)\kappa_{\rm Carr}\). The laws
\eqref{eq:ley-carry-correcta-ch} and
\eqref{eq:carry-no-parametrico-ch}
give \eqref{eq:carry-injerto-ch}. The ledger preserves
\((\mathfrak p_\varepsilon,2\varepsilon,r_3(2\varepsilon),
\chi(\varepsilon,\varepsilon))\); the three values of this ledger are
distinct even when the scalar increment is zero. Hence the images are not
identified, and the block decoder recovers a unique \(\varepsilon\).

The edge definitions depend only on the new phase and concatenate, rather than
replace, the prefix, cylinder, boundary, and ledger. They are also invariant
under
\(\jmath_d\operatorname{mem}_d\mapsto
  \jmath_d\operatorname{mem}_d+nu\)
and under translation of the input carry. Consequently, the same construction
applies at the endpoint of any turn. For every finite state, the explicit
sequence
\[
 x,\quad\gamma_{0,k_n}(x),\quad
 \gamma_{0,k_{n+1}}\gamma_{0,k_n}(x),\quad\ldots
\]
is a compatible infinite tail. This proves the survival of all constructed
states and the repeated availability of the three branches, without
postulating either property. Consequently, in accordance with the preceding
level definitions, we may finally identify
\begin{equation}
 H_{k_n}^{\rm cal}:=G_n
 \label{eq:subarbol-calibrado-superviviente-ch}
\end{equation}
for every \(n\ge0\).
\end{proof}

The preceding closure distinguishes the three clocks involved. Nine emissions
return the phase in \(C_9\) and advance memory. Twelve turns, that is,
\(108\) emissions, also return the position of the observable calendar; even
then, neither the ledger nor memory is erased. The operator
\(\mathcal K_{\rm ph}\) acts only afterward, as a recognition of reduced
memory from the already constructed turns; it never generates the edges of
\(\operatorname{Ext}^{\rm enr,cal}\).
